\documentclass[a4paper,10pt]{article}
\usepackage{lamia-pedagogy}

\begin{document}

% =========================================================================
% PAGE 1 : Fiche Enseignant & Démarche Pédagogique
% =========================================================================
\ficheHeader{1 Heure}{\makebox[2.5cm]{\dotfill}}{1\textsuperscript{ère} AM}{2026 / 2027}

\vspace{-3pt}
\competenceBanner{\faHeadphones\ Projet 01 / Séquence 01 — Séance 01 : Compréhension de l'Oral}

\vspace{-2pt}
\modaliteSupportBar{Collectif puis individuel}{Support audio / Script transcrit \guillemotleft\ Le lavage des mains selon M. Pompe \guillemotright}

\vspace{1pt}

\begin{minipage}[t]{0.485\textwidth}
\begin{tcolorbox}[
  enhanced,
  colback=secondarylight,
  colframe=secondary,
  arc=1.5mm,leftrule=3pt,rightrule=0.5pt,toprule=0.5pt,bottomrule=0.5pt,
  top=3pt,bottom=3pt,left=6pt,right=6pt,
  title={\small\textbf{\faBullseye\ Objectifs d'apprentissage}},
  coltitle=secondary!90!black,attach boxed title to top left={yshift=-2mm,xshift=4mm},
  boxed title style={colback=white,boxrule=0.5pt,colframe=secondary,arc=1mm,top=1pt,bottom=1pt}
]
\vspace{1pt}
{\small
\begin{itemize}[leftmargin=*,noitemsep,topsep=1pt]
  \item[\faCheckCircle] Identifier le thème global du message oral.
  \item[\faCheckCircle] Retrouver des informations précises à travers 3 écoutes.
  \item[\faCheckCircle] Comprendre le rôle du savon contre les microbes.
\end{itemize}
}
\end{tcolorbox}
\end{minipage}%
\hfill
\begin{minipage}[t]{0.485\textwidth}
\begin{tcolorbox}[
  enhanced,
  colback=accentlight!35,
  colframe=accent,
  arc=1.5mm,leftrule=3pt,rightrule=0.5pt,toprule=0.5pt,bottomrule=0.5pt,
  top=3pt,bottom=3pt,left=6pt,right=6pt,
  title={\small\textbf{\faLightbulb\ Éveil de l'intérêt}},
  coltitle=accent!90!black,attach boxed title to top left={yshift=-2mm,xshift=4mm},
  boxed title style={colback=white,boxrule=0.5pt,colframe=accent,arc=1mm,top=1pt,bottom=1pt}
]
\vspace{1pt}
{\small
\textbf{Questions d'amorce :}
\begin{itemize}[leftmargin=*,noitemsep,topsep=1pt]
  \item Pourquoi se lave-t-on les mains ?
  \item Quand doit-on se laver les mains ?
  \item Que risque-t-on si nos mains sont sales ?
\end{itemize}
}
\end{tcolorbox}
\end{minipage}

\vspace{3pt}

% Transcription du Support
\begin{tcolorbox}[
  enhanced,
  colback=white,
  colframe=primary!70,
  arc=2mm,boxrule=0.8pt,
  top=5pt,bottom=5pt,left=8pt,right=8pt,
  title={\small\textbf{\textcolor{primary}{\faVolumeUp\ Texte du support d'écoute : \guillemotleft\ Le lavage des mains selon M. Pompe \guillemotright}}},
  coltitle=primary,colbacktitle=primarylight,
  titlerule=0.5pt,titlerule style=primary!40
]
{\small
\guillemotleft\ Bonjour les enfants ! Je m'appelle Monsieur Pompe, le distributeur de savon liquide de votre école. Aujourd'hui, je viens vous parler d'un geste tout simple mais très important : \textbf{le lavage des mains}.

Sur vos mains, il y a des millions de petits êtres vivants invisibles à l'\oe il nu qu'on appelle \textbf{les microbes}. Quand vous touchez des objets, la terre ou les poignées de porte, les microbes s'accrochent à votre peau. Si vous mangez avec des mains sales, ces microbes entrent dans votre ventre et vous rendent malades !

Pour les éliminer, l'eau seule ne suffit pas : il faut toujours m'utiliser avec de l'eau propre. Frottez bien vos paumes, le dos des mains et entre les doigts pendant au moins 30 secondes, puis rincez et séchez-vous avec une serviette propre. Des mains propres, c'est la garantie d'une bonne santé ! \guillemotright
}
\end{tcolorbox}

\vspace{3pt}

% Démarche pédagogique
\begin{tcolorbox}[
  enhanced,
  colback=neutrallight,
  colframe=cardborder,
  arc=1.5mm,boxrule=0.6pt,
  top=4pt,bottom=4pt,left=7pt,right=7pt,
  title={\small\textbf{\textcolor{primary}{\faChalkboardTeacher\ Démarche des 3 écoutes méthodiques}}},
  coltitle=primary,colbacktitle=white,titlerule=0.4pt
]
{\footnotesize
\begin{tabularx}{\linewidth}{l X}
\textbf{1\textsuperscript{ère} écoute :} & Découverte globale sans prise de notes. Compréhension du thème et des personnages.\\
\textbf{2\textsuperscript{ème} écoute :} & Écoute sélective. Repérage des mots clés : \textit{microbes}, \textit{savon}, \textit{maladies}, \textit{30 secondes}.\\
\textbf{3\textsuperscript{ème} écoute :} & Écoute fine et vérification. Réponse aux questions de synthèse et restitution orale.\\
\end{tabularx}
}
\end{tcolorbox}

\newpage

% =========================================================================
% PAGE 2 : Fiche Élève (Activités d'Écoute et de Compréhension)
% =========================================================================
\ficheEleveHeader

\vspace{-3pt}
\competenceBanner{\faPenFancy\ Fiche d'Activités : Compréhension de l'Oral}

\vspace{2pt}

% ACTIVITÉ 1 : 1ère Écoute
\begin{tcolorbox}[
  enhanced,
  colback=white,
  colframe=cardborder,
  arc=2mm,boxrule=0.8pt,
  top=4pt,bottom=4pt,left=8pt,right=8pt,
  title={\small\textbf{\textcolor{primary}{\faPlay\ 1\textsuperscript{ère} Écoute : Compréhension globale}} \hfill \badge[primary]{Écoute 1}},
  coltitle=primary,colbacktitle=primarylight,titlerule=0.5pt,titlerule style=cardborder
]
{\small
\textbf{Consigne :} Coche la bonne réponse.
\begin{enumerate}[label=\textbf{\arabic*.},leftmargin=*,noitemsep,topsep=2pt]
  \item Qui parle dans ce document sonore ?
    \cardOptThree{Un médecin}{Monsieur Pompe (le savon)}{Un élève}
  \item Quel est le sujet principal du document ?
    \cardOptThree{La pratique du sport}{Le lavage des mains}{Les vacances}
\end{enumerate}
}
\end{tcolorbox}

\vspace{3pt}

% ACTIVITÉ 2 : 2ème Écoute
\begin{tcolorbox}[
  enhanced,
  colback=white,
  colframe=cardborder,
  arc=2mm,boxrule=0.8pt,
  top=4pt,bottom=4pt,left=8pt,right=8pt,
  title={\small\textbf{\textcolor{primary}{\faStepForward\ 2\textsuperscript{ème} Écoute : Analyse et repérage}} \hfill \badge[secondary]{Écoute 2}},
  coltitle=primary,colbacktitle=primarylight,titlerule=0.5pt,titlerule style=cardborder
]
{\small
\textbf{Consigne :} Réponds par \textbf{Vrai} ou \textbf{Faux}.
\begin{enumerate}[label=\textbf{\alph*)},leftmargin=*,noitemsep,topsep=2pt]
  \item Les microbes sont visibles facilement à l'\oe il nu. \dotfill \cardOptTwo{Vrai}{Faux}
  \item L'eau seule suffit pour éliminer tous les microbes. \dotfill \cardOptTwo{Vrai}{Faux}
  \item Il faut se frotter les mains pendant au moins 30 secondes. \dotfill \cardOptTwo{Vrai}{Faux}
  \item Les microbes peuvent causer des maladies graves. \dotfill \cardOptTwo{Vrai}{Faux}
\end{enumerate}
}
\end{tcolorbox}

\vspace{3pt}

% ACTIVITÉ 3 : 3ème Écoute & Synthèse
\begin{tcolorbox}[
  enhanced,
  colback=white,
  colframe=cardborder,
  arc=2mm,boxrule=0.8pt,
  top=4pt,bottom=4pt,left=8pt,right=8pt,
  title={\small\textbf{\textcolor{primary}{\faCheckDouble\ 3\textsuperscript{ème} Écoute : Récapitulation}} \hfill \badge[emerald]{Synthèse}},
  coltitle=primary,colbacktitle=primarylight,titlerule=0.5pt,titlerule style=cardborder
]
{\small
\textbf{Consigne :} Complète le paragraphe suivant à l'aide des mots : \textbf{savon}, \textbf{malades}, \textbf{mains}, \textbf{microbes}.

\vspace{4pt}
Se laver les \drule{2.8cm} avec de l'eau et du \drule{2.8cm} permet d'éliminer les \drule{2.8cm} et d'éviter de tomber \drule{2.8cm}.
}
\end{tcolorbox}

\vspace{3pt}

% Grille d'auto-évaluation
\begin{tcolorbox}[
  enhanced,
  colback=neutrallight,
  colframe=secondary!60!cardborder,
  arc=1.5mm,boxrule=0.6pt,
  top=3pt,bottom=3pt,left=7pt,right=7pt,
  title={\scriptsize\textbf{\faTasks\ Auto-évaluation de l'élève}},
  coltitle=primary,colbacktitle=secondarylight,titlerule=0.4pt
]
{\footnotesize
\begin{tabularx}{\linewidth}{Xcc}
\textbf{Je suis capable de :} & \textbf{Oui \faSmile} & \textbf{Non \faFrown} \\
\midrule
Identifier qui parle et le thème abordé. & \qcmCircle & \qcmCircle \\
Donner la durée recommandée pour un bon lavage de mains. & \qcmCircle & \qcmCircle \\
\end{tabularx}
}
\end{tcolorbox}

\end{document}
